\documentclass{article}
\usepackage{amsmath,amssymb}
\begin{document}
The problem asks for a computably represented countable structure that
is computably AUT-countable on a cone, but for which no finite parameter
set gives $\Sigma^{in}_1$ definitions of all automorphisms from the
parameters and their images.  The construction below actually gives the
stronger property that every automorphism of every copy is computable
from that copy, with no cone oracle.

Working construction:
Let $V=[\omega]^{<\omega}$ be the countable $\mathbb F_2$-vector space
under symmetric difference.  Fix a computable sequence
$(D_i)_{i\in\omega}$ of finite subsets of $V$ in which every finite
subset occurs cofinally often.  Code the two-sorted structure as the
one-sorted computable disjoint union
$(\{0\}\times\omega)\sqcup(\{1\}\times V)$, with unary predicates $I$ and
$S$.  Interpret $<$ as the usual order on $I$ and false outside
$I^2$.  Interpret $R(i,x,y)$, for $i\in I$ and $x,y\in S$, by
$x+y\in D_i$, and make $R$ false on all other triples.

Automorphism group:
The ordered sort $I\cong(\omega,<)$ is rigid, hence fixed pointwise by
every automorphism.  Since singleton sets $\{h\}$ occur among the
$D_i$, preservation of the corresponding marker gives
$f(x+h)=f(x)+h$ for all $x,h\in V$.  Thus $f(h)=f(0)+h$, so all
automorphisms are translations $\tau_c:x\mapsto x+c$ on the $S$-sort,
fixing $I$ pointwise; conversely all translations preserve differences
and hence preserve $R$.

Computability in arbitrary copies:
If $\mathcal B\cong\mathcal A$ and $\sigma$ transports to $\tau_c$, pick
in $\mathcal B$ a marker corresponding to some standard $i$ with
$D_i=\{c\}$.  A machine with this marker hardwired and oracle the atomic
diagram of $\mathcal B$ fixes the $I$-sort and, on an $S$-input $x$,
searches for the unique $y$ with $R(i,x,y)$.  This is a non-uniform
relative computation of the particular automorphism, as required by the
definition.

Nondefinability argument:
Given finite parameters $P$, put $M=\max(P\cap I)$ or $-1$, and
$W=\operatorname{span}\bigcup_{i\le M}D_i$.  Then $W$ is finite.  Choose
nonzero $h\notin W$, let $\pi=\tau_h$, set $Q=P\cup\pi(P)$, and choose
$x$ so the cosets $x+W$ and $x+h+W$ avoid $Q\cap S$.  Let $y=x+h$ and
$C=y+W$.

For any existential formula over $Q$ true of $(x,y)$, choose witnesses
and a true conjunction of literals in disjunctive normal form.  Let $F$
be the finite set of $S$-sort values appearing.  Shift exactly the
values in $F\cap C$ by a fresh vector $t$.  Avoid finitely many $t$'s:
$t=0$; collision values between shifted and unshifted $S$-values; values
that would make a cross-cut difference land in a low $D_i$ for
$i\le M$; and values that would make, for the same high marker equality
class, a positive high-marker difference equal a negative high-marker
difference.  The latter obstruction is finite because each resulting
difference is either its old value or old value plus $t$, so equality is
either already impossible from the original true diagram or imposes one
linear equation on $t$.

After choosing such a $t$, low-marker $R$-literals are preserved.  For
each high marker equality class $\mu$, collect the new positive
differences $A_\mu$ and negative differences $B_\mu$; they are disjoint.
Replace the high marker witnesses by fresh indices above $M$, preserving
their finite equality/order pattern, and satisfying $D_i=A_\mu$ for the
corresponding class.  Cofinal repetition of every finite subset among
the $D_i$ makes the simultaneous choice possible.  Sort literals,
equality, order, wrong-sort $R$ atoms, low $R$ atoms, and high $R$ atoms
are all preserved.  The same formula holds of $(x,y+t)$, but
$y+t\ne\pi(x)$, so no existential formula true of $(x,\pi(x))$ can be
contained in the graph of $\pi$.  Hence no $\Sigma^{in}_1$ definition
over $P\cup\pi(P)$ defines $\pi$.
\end{document}
